\chapter{Mass Storage: SD Cards and the ext4 Filesystem}
\label{ch:storage}

The final layer is mass storage: two block drivers for SD cards, and a complete
pure-Ada ext2/3/4 filesystem (a reimplementation of lwext4) that sits on top of
them and mounts real \texttt{mkfs.ext4} cards.

\section{SD\_SPI --- an SD card over the SPI master}

\texttt{ESP32S3.SD\_SPI} talks to an SD/SDHC card in SPI mode, layered on the
task-safe SPI driver. The chip-select is an ordinary GPIO so it can be held across
a whole command. Blocks are 512-byte logical sectors (LBA addressing); SDSC and
SDHC are detected and handled internally.

\begin{lstlisting}
type Block is array (0 .. 511) of Interfaces.Unsigned_8;
type Status is (OK, No_Card, Unusable, Init_Timeout, Read_Error, Write_Error);
procedure Setup (C : out Card; Host : ESP32S3.SPI.SPI_Host;
                 Sclk, Mosi, Miso, Cs : Pin_Id;
                 Init_Clock_Hz : Positive := 400_000;
                 Data_Clock_Hz : Positive := 8_000_000);
procedure Initialize (C : in out Card; Result : out Status);
procedure Read_Block  (C : in out Card; LBA : Block_Address;
                       Data : out Block; Result : out Status);
procedure Write_Block (C : in out Card; LBA : Block_Address;
                       Data : Block; Result : out Status);
\end{lstlisting}

\textbf{Example 1 --- bring up a card and read sector 0.}
\begin{lstlisting}
with ESP32S3.SD_SPI; use ESP32S3.SD_SPI;
with ESP32S3.SPI;
procedure SD_Init is
   C   : Card;
   St  : Status;
   Sec : Block;
begin
   Setup (C, ESP32S3.SPI.SPI2, Sclk => 12, Mosi => 11, Miso => 13, Cs => 10);
   Initialize (C, St);
   if St = OK then
      Read_Block (C, LBA => 0, Data => Sec, Result => St);   --  the MBR sector
   end if;
end SD_Init;
\end{lstlisting}

\textbf{Example 2 --- a non-destructive read / write-back / re-read round-trip.}
\begin{lstlisting}
with ESP32S3.SD_SPI; use ESP32S3.SD_SPI;
with ESP32S3.SPI;
procedure SD_Roundtrip is
   C          : Card;
   St         : Status;
   Orig, Back : Block;
begin
   Setup (C, ESP32S3.SPI.SPI2, Sclk => 12, Mosi => 11, Miso => 13, Cs => 10);
   Initialize (C, St);
   Read_Block  (C, 16#2000#, Orig, St);     --  read a scratch sector
   Write_Block (C, 16#2000#, Orig, St);     --  write the same bytes back
   Read_Block  (C, 16#2000#, Back, St);     --  re-read; Back = Orig, no data lost
end SD_Roundtrip;
\end{lstlisting}

\section{SDMMC --- the native SD/MMC host}

\texttt{ESP32S3.SDMMC} drives the dedicated SDHOST controller (clock + command +
1- or 4-bit data) in PIO mode --- faster than SPI mode, and with no DMA or
finalization it works under every profile. The API mirrors \texttt{SD\_SPI}.

\begin{lstlisting}
procedure Setup (C : out Card; On : Slot; Clk, Cmd, D0 : Pin_Id;
                 D1, D2, D3 : Optional_Pin := No_Pin;
                 Width : Bus_Width := Width_1; ...);
procedure Initialize (C : in out Card; Result : out Status);
procedure Read_Block  (C : in out Card; LBA : Block_Address; Data : out Block;
                       Result : out Status);
\end{lstlisting}

\textbf{Example --- 4-bit bus, read a sector.}
\begin{lstlisting}
with ESP32S3.SDMMC; use ESP32S3.SDMMC;
procedure SDMMC_Read is
   C   : Card;
   St  : Status;
   Sec : Block;
begin
   Setup (C, On => Slot1, Clk => 14, Cmd => 15, D0 => 2, D1 => 4, D2 => 12,
          D3 => 13, Width => Width_4);
   Initialize (C, St);
   if St = OK then
      Read_Block (C, 0, Sec, St);
   end if;
end SDMMC_Read;
\end{lstlisting}

\section{The ext4 filesystem}

\texttt{ESP32S3.Ext4} is a pure-Ada ext2/3/4 implementation with a JBD2 journal
(Chapter overview in \texttt{libs/esp32s3\_hal/src/ext4/README.md}). It reads
default \texttt{mkfs.ext4} cards (extents, 64-bit, HTree directories,
\texttt{metadata\_csum} verification), writes ext2/3/non-csum filesystems, and
recovers a dirty journal on mount. It needs the \texttt{embedded} or \texttt{full}
profile.

\subsection{Connecting a block device}

The filesystem talks to a small block-device abstraction (\texttt{ESP32S3.%
Block\_Dev.Device}). An adapter turns an initialised SD card into one with a
single call --- the \texttt{Make} below appears in each complete example that
follows:

\begin{lstlisting}
with ESP32S3.Block_Dev;
with ESP32S3.Block_Dev.SD_SPI_Source;
--  ... given an initialised "Card : aliased ESP32S3.SD_SPI.Card" ...
   Dev : constant ESP32S3.Block_Dev.Device :=
           ESP32S3.Block_Dev.SD_SPI_Source.Make (Card'Access);
\end{lstlisting}

\subsection{Mounting and reading}

A \texttt{Mount} is a controlled handle: it flushes and closes on scope exit, and
\texttt{Open} automatically replays a dirty journal on a writable mount. The
complete example brings up the card, wraps it, mounts read-only, then reads a
file:

\begin{lstlisting}
with ESP32S3.SPI;
with ESP32S3.SD_SPI;
with ESP32S3.Block_Dev;
with ESP32S3.Block_Dev.SD_SPI_Source;
with ESP32S3.Ext4;
with ESP32S3.Ext4.FS;
with ESP32S3.Ext4.Inode;
procedure EXT4_Read is
   use type ESP32S3.SD_SPI.Status;
   Card : aliased ESP32S3.SD_SPI.Card;
   St   : ESP32S3.SD_SPI.Status;
begin
   ESP32S3.SD_SPI.Setup (Card, ESP32S3.SPI.SPI2,
                         Sclk => 12, Mosi => 11, Miso => 13, Cs => 10);
   ESP32S3.SD_SPI.Initialize (Card, St);
   if St /= ESP32S3.SD_SPI.OK then
      return;
   end if;
   declare
      Dev  : constant ESP32S3.Block_Dev.Device :=
               ESP32S3.Block_Dev.SD_SPI_Source.Make (Card'Access);
      M    : ESP32S3.Ext4.FS.Mount;
      I    : ESP32S3.Ext4.Inode.Info;
      Buf  : ESP32S3.Ext4.Byte_Array (0 .. 255);
      Last : Natural;
   begin
      M.Open (Dev, Read_Only => True);
      M.Stat (M.Lookup ("/etc/hello.txt"), I);  --  resolve a path -> its inode
      M.Read_File (I, Offset => 0, Into => Buf, Last => Last);
      --  Buf (0 .. Last - 1) holds the first Last bytes of the file
   end;   --  M finalizes: flush + close
end EXT4_Read;
\end{lstlisting}

\textbf{Listing a directory.} \texttt{Iterate} calls a visitor for each entry.
The card bring-up is identical to the example above; the body is:

\begin{lstlisting}
   declare
      Dev  : constant ESP32S3.Block_Dev.Device :=
               ESP32S3.Block_Dev.SD_SPI_Source.Make (Card'Access);
      M    : ESP32S3.Ext4.FS.Mount;
      Root : ESP32S3.Ext4.Inode.Info;
      procedure Show (Name : String;
                      Ino  : ESP32S3.Ext4.Inode_Number;
                      File_Type : ESP32S3.Ext4.U8) is
      begin
         Put_Line ("  " & Name);
      end Show;
   begin
      M.Open (Dev, Read_Only => True);
      M.Stat (M.Lookup ("/"), Root);
      M.Iterate (Root, Show'Access);
   end;
\end{lstlisting}

\subsection{Writing}

On a non-\texttt{metadata\_csum} filesystem (ext2/3, or
\texttt{mke2fs -O \textasciicircum metadata\_csum}) the filesystem can be modified.
The full POSIX-like set is available: \texttt{Create\_File}, \texttt{Write\_File},
\texttt{Truncate}, \texttt{Mkdir}, \texttt{Rmdir}, \texttt{Unlink},
\texttt{Rename} and hard \texttt{Link}.

\begin{lstlisting}
   declare
      Dev : constant ESP32S3.Block_Dev.Device :=
              ESP32S3.Block_Dev.SD_SPI_Source.Make (Card'Access);
      M   : ESP32S3.Ext4.FS.Mount;
      N   : ESP32S3.Ext4.Inode_Number;
   begin
      M.Open (Dev, Read_Only => False);
      N := M.Create_File ("/", "log.txt");
      M.Write_File (N, (Character'Pos ('h'), Character'Pos ('i'), 10));
      M.Mkdir ("/", "data");
      M.Close;          --  flush + close (or let M finalize)
   end;
\end{lstlisting}

\subsection{Crash-safe transactions}

For atomicity, group an operation behind \texttt{Commit}: the dirty metadata and
the superblock are written to the journal (durably), the RECOVER flag is set ---
the commit barrier --- then the metadata is checkpointed to its final locations
and the flag cleared. A power loss after the barrier is recovered forward by the
next \texttt{Open} (or by any other ext4 implementation, including the Linux
kernel); a loss before it has no effect.

\begin{lstlisting}
   --  Payload : constant ESP32S3.Ext4.Byte_Array := (16#41#, 16#42#, 16#43#);
   N := M.Create_File ("/", "important.txt");
   M.Write_File (N, Payload);
   M.Commit;         --  one atomic, journaled, crash-safe transaction
\end{lstlisting}

\subsection{Errors}

Filesystem operations raise: the standard \texttt{Ada.IO\_Exceptions} names
(\texttt{Name\_Error} for a missing path, \texttt{Use\_Error}, \texttt{Device\_%
Error}, \texttt{End\_Error}) plus filesystem-specific \texttt{Unsupported\_%
Feature}, \texttt{Bad\_Checksum}, \texttt{Corrupt}, \texttt{No\_Space},
\texttt{Not\_Empty} and \texttt{Read\_Only}. Wrap calls in a handler so a bad card
or a missing file never escapes a task:

\begin{lstlisting}
begin
   M.Stat (M.Lookup ("/maybe/missing"), I);
exception
   when ESP32S3.Ext4.Name_Error => Put_Line ("not found");
   when ESP32S3.Ext4.Bad_Checksum => Put_Line ("corrupt metadata");
end;
\end{lstlisting}

\subsection{How it is validated}

Because the filesystem is pure logic over the block interface, it is developed and
tested host-native (x86) against real \texttt{mke2fs} images, with the Linux
kernel and \texttt{e2fsck} as the oracle: every read is byte-compared to the
source file, every mutation is checked with \texttt{e2fsck -fn}, and journal
replay and commit are cross-checked against the kernel's own recovery. The same
code cross-compiles for the ESP32-S3 and runs over the SD drivers above. See
\texttt{tests/ext4\_host/} and \texttt{libs/esp32s3\_hal/src/ext4/REFERENCES.md}.

\vfill
\noindent\rule{\linewidth}{0.4pt}\\[4pt]
\textit{The whole stack --- runtime, peripheral HAL, and filesystem --- is one
coherent body of Ada: strong types from the register up, the same RAII and
task-safety patterns in every driver, and a filesystem validated against the
Linux kernel, all running on both cores of the ESP32-S3 on its own bare-metal
runtime.}
